% Source: https://artofproblemsolving.com/community/c3h2826137_2022_abmc_team_rounds_58__actonboxborough_math_competition
% Stable identity: aops-topic-2826137-q13
The hexagon apothem is $1$, so its circumradius and side length are $\dfrac{2}{\sqrt3}$. Its area is half the perimeter times the apothem, or $2\sqrt3$, while the outer circle has area $\dfrac{4\pi}{3}$. Thus $(a,b,c,d)=(4,3,2,3)$ and the answer is $\boxed{12}$.
\begin{asy}
usepackage("amsmath");
size(7cm);
real R=2/sqrt(3);draw(circle((0,0),R));draw(circle((0,0),1),gray);path h=R*dir(0);for(int k=1;k<6;++k)h=h--R*dir(60*k);draw(h--cycle,blue);draw((0,0)--dir(30));label("$1$",.5*dir(30),NW);
\end{asy}
